\section{Problem Definition}
\label{sec:problem_definition}

This paper studies the problem of scheduling driver activities for long-distance freight transport with a battery electric truck (BET). The schedule must simultaneously satisfy Hours-of-Service (HoS) regulations and battery energy constraints. Given a predefined sequence of customer visits, the objective is to minimize the total route duration by determining the optimal timing and locations of mandatory breaks, daily rest periods, and battery charging operations.

\begin{anew}
\textbf{Missing: a short assumptions list.} Suggested as a new \S3.1, before
``Route and Customers''. Everything in it is already true of the code; two items
(no idle waiting, deterministic queue) do real work in Section~\ref{sec:diesel}
and are currently discoverable only by reading the source.

\medskip
\noindent\textbf{3.1 Modelling assumptions.}
(A1) The order of customer visits is given; routing is out of scope.
(A2) One truck executes one route between two weekly rests; fleet effects and
depot slow charging are excluded.
(A3) There is no idle waiting: service starts on arrival and the only way to
delay a departure is to declare a break or a rest, so an early arrival is
penalised rather than absorbed.
(A4) Only the daily provisions of Reg.~561/2006 and Dir.~2002/15/EC are modelled;
weekly rests, the 56\,h weekly driving cap and the 60\,h working week are out of
scope, and $\bar\rho$ and $\bar\epsilon$ are budgeted per route, which is exact
when the route is completed between two weekly rests.
(A5) Energy follows a speed-dependent rate evaluated at the realized leg-average
speed; grade, payload variation and the extra consumption of stop-and-go traffic
at a given average speed are not represented.
(A6) Charger access delay $Q_i$ is a known instance parameter, identical for
every method; charger availability and reservation are out of scope.
(A7) The objective is makespan plus a fixed service-level penalty; energy price,
driver wage and charging tariffs are not priced.
\end{anew}

\subsection{Route and Customers}
The route connects a sequence of customer locations via road segments whose lengths are known in advance. Travel on each segment is characterized by an average speed. In practice, however, traffic conditions introduce uncertainty in the actual travel duration: this is explicitly accounted for in the model. The total route duration is the sum of five components: driving time accumulated while the vehicle is moving between locations, service time for loading or unloading operations, charging time at charging stations, station access and manoeuvering time, and legally required interruptions such as mandatory breaks and daily rest periods.

The sequence of customers to be visited is fixed and known prior to departure. At each customer location, the driver performs a loading or unloading operation whose duration is known.
At each customer, soft time windows are considered. Service begins immediately upon arrival, whether or not the window is open: distribution terminals do not turn a truck away, but an arrival outside the announced window, either early or late, disrupts terminal planning and incurs a fixed service-level penalty, identical for early and late arrivals. Any break or rest taken at a customer stop occurs after service. \NEW{The vehicle is not permitted to wait idly for a window to open: an early arrival is served early and penalised, or the driver stretches a break (assumption A3).}

\begin{atodo}
\textbf{Justify $\beta$.} The paper's central uncertainty finding is that
uncertainty is paid in windows rather than in driving time, and the size of that
finding is proportional to $\beta$, which Table~\ref{tab:params} sets to 1\,h with
no argument. \texttt{settings.py} already carries
\texttt{BETA\_TW\_CLASSES = (1, 2, 5)}. Either run that sweep on the reduced set,
or say in one sentence why one hour of makespan is the right price for one missed
appointment. The duration-only gaps already reported in parentheses in
Table~\ref{tab:gap} are your insurance either way.
\end{atodo}

\subsection{Charging operations}
The BET has a battery of fixed capacity (kWh). Energy consumption is assumed to be dependent of the distance traveled and the vehicle speed, yielding a specific consumption rate expressed in kWh/km. At each CS stop, the driver may charge the battery either partially or fully, and multiple charging stops are permitted throughout the route. Charging is not allowed during customer service operations.

Only high-power fast charging is available along the route (overnight slow charging such as in depot is not considered). The charging process is modeled by a piecewise linear approximation of the nonlinear charge-rate curve. Queuing at a CS is modeled as a known waiting time at each CS\FIX{, drawn once at instance generation and visible to every method, so it is a deterministic cost of using a station rather than a source of uncertainty}.

Finally, a sufficiently long charging stop may simultaneously satisfy a mandatory break requirement. Specifically, a charging event counts as a driving break provided its duration meets the applicable threshold: at least 15~minutes for a first split break, at least 30~minutes for a second split break, or at least 45~minutes for a full break, as defined by the HoS regulations described below.

\begin{atodo}
Give the legal basis for ``charging counts as a break'' in a footnote: Art.~7 of
Reg.~561/2006 requires only that the driver not drive and be free to dispose of
the time, and supervising a charge that requires no intervention is generally
read as compatible. This is the load-bearing assumption of the whole paper and it
is currently asserted without authority.
\end{atodo}

\subsection{HoS Regulations}
The scheduling problem is subject to European HoS regulations given by Regulation (EC) No 561/2006 (driving times, breaks and rests) and Directive 2002/15/EC (working time). \FIX{We model the daily provisions of both instruments; weekly-level accounting is outside the scope of this study, and the reduced-rest and extended-driving allowances are consequently budgeted per route (assumption A4).}
Two accumulation counters are maintained: driving time, which counts only periods when the vehicle is in motion, and working time, which additionally includes customer service, charger access and maneuvering, and charging time not covered by a simultaneous break.

Regarding mandatory breaks, a driver may not accumulate more than 4.5~hours of continuous driving without taking a full break of at least 45~minutes. This full break may be replaced by a split sequence consisting of a first period of at least 15~minutes (but less than 45~minutes) followed, in that order, by a second period of at least 30~minutes.

Regarding shift-level constraints, the maximum permitted driving time per shift is 9~hours, extendable to 10~hours on no more than two occasions per week.
A daily rest of at least 11 h must separate consecutive shifts, and it may be reduced to 9 h at most three times between weekly rests.
As a consequence of these rest requirements, the sum of working time, break time, and any other on-duty periods within a 24-hour window is bounded by 13~hours under a full rest, or by 15~hours when the rest is reduced to 9~hours.

Charging during rests is treated as sequential: the truck must vacate the charging bay before the rest begins, and such charging time counts as work.

While charging is only possible at charging stations,  breaks and rests may be taken at customer stops, charging stations, and dedicated rest-area nodes inserted along legs.

\section{Deterministic Mathematical Model} \label{sec:model}
We now present the mathematical formulation for the deterministic problem. Tables ~\ref{tab:notation-mini} and ~\ref{tab:notation-mini2} present all notation used in the model formulation. Figure \ref{fig:solution} present a possible solution of the problem.

\begin{table}[h!]
\centering
\caption{Notation summary.}
\label{tab:notation-mini}
\small
\begin{tabular}{p{0.12\linewidth}p{0.8\linewidth}}
\toprule
\multicolumn{2}{l}{\textbf{Sets}}\\
\midrule
$\mathcal{I}$ & Set of possible stops, ordered from 0=origin to N=destination.\\
$\mathcal{C}$ & Set of customers $\mathcal C \subseteq \mathcal I$.\\
$\mathcal{K}$ & Set of charging stations $\mathcal K \subseteq \mathcal I$.\\
$\mathcal{L}$ & Set of rest areas $\mathcal L \subseteq \mathcal I$, $\mathcal C \cap \mathcal K \cap \mathcal L= \emptyset$.\\
$\mathcal{R}$ & Set of breakpoints of the piecewise-linear charging function.\\
\midrule
\multicolumn{2}{l}{\textbf{Parameters}}\\
\midrule
$D_i$ & Travel time from stop $i$ to stop $i+1$.\\
$E_i$ & Energy consumption between stop $i$ and stop $i+1$.\\
$E^{cap}$ & Battery capacity. \\
$E^{\min}$ & Minimum SOC. \\
$E^0$ & Initial SOC. \\
$\bar{E}_r$  & SOC level at breakpoint $r \in \mathcal{R}$ of the charging fct, ($\bar{E}_0 = 0$, $\bar{E}_R = E^{cap}$).\\
$\mathcal{T}_r$ & Time to charge from empty to $\bar{E}_r$,($\mathcal{T}_0 = 0$).\\
$\mathcal{T}_{\text{min}}$ & Minimum charging time at a CS \FIX{($= 15$\,min)}.\\
$S_i$ & Service time at customer $i \in \mathcal{C}$. \\
$Q_i$ & Queue time at CS stop $i \in \mathcal{K}$ if charging occurs. Zero if $y_i = 0$.\\
$M_i^{stop}$ & Maneuver time at CS stop $i \in \mathcal{S}$ if any activity occurs.\\
$M_i^{seq}$ & Repositioning time at stop $i \in \mathcal{K}$ when charging and a break or rest are declared sequentially, i.e. the truck must vacate the charging bay before the break begins.\\
$W_i^{Ha}, W_i^{Hf}$  & Service time window at customer $i \in \mathcal{C}$. \\
$ \beta$ & Fixed out-of-window service penalty (h). \\
$T^{brk}_{b45}$   & Full break duration ($= 45$\,min). \\
$T^{brk}_{b_{15}}$     & Min first split break ($= 15$\,min). \\
$T^{brk}_{b_{30}}$     & Min second split break ($= 30$\,min). \\
$T^{rst}_{r_1}$          & Minimum daily rest duration ($= 11$\,h). \\
$T^{rst}_{r_2}$          & Reduced daily rest duration (exception 1) ($= 9$\,h at most $3\times$/week). \\
$\bar{\rho}$ & \DEL{Weekly n}\FIX{N}umber of allowed reduced rest\FIX{s per route} ($=3$).\\
$T^{drv}_{cons}$   & Max consecutive driving time ($= 4.5$\,h). \\
$T^{drv}_{shift_1}$     & Max shift driving time ($= 9$\,h). \\
$T^{drv}_{shift_2}$     & Increased shift driving time (exception 2) ($= 10$\,h at most $2\times$/week). \\
$\bar \epsilon $ &  \DEL{Weekly n}\FIX{N}umber of allowed extended driving shifts \FIX{per route} ($=2$).\\
\DEL{$T^{wk}_{sh1}$}\FIX{$T^{spr}_{1}$}      & Max shift \DEL{working time}\FIX{spread} under a regular daily rest ($= 13$\,h). \\
\DEL{$T^{wk}_{sh2}$}\FIX{$T^{spr}_{2}$}      & Max shift \DEL{working time}\FIX{spread} under a reduced daily rest ($= 15$\,h). \\
\bottomrule
\multicolumn{2}{l}{{\footnotesize Abbreviations: SOC = state of charge, CS = charging station}}\\
\end{tabular}
\end{table}

\begin{acmt}
Three notation fixes. (i) $\mathcal T_{\min}$ was listed but never given a value
anywhere; \texttt{MILP.py} fixes it at 15\,min
(\texttt{tauc[i] >= 0.25*y[i]}). (ii) $T^{wk}_{sh1}, T^{wk}_{sh2}$ here are the
same objects the constraint section calls $T^{spr}_1, T^{spr}_2$, and they bound
the \emph{spread}, not the working time --- two names and two meanings for one
parameter. (iii) $\bar\rho$ and $\bar\epsilon$ are described as weekly but
enforced once per route in \eqref{eq:r2_lim} and \eqref{eq:ext_budget}.
\end{acmt}

\begin{table}[h!]
\centering
\caption{Variables notation summary.}
\label{tab:notation-mini2}
\small
\begin{tabular}{p{0.25\linewidth}p{0.75\linewidth}}
\toprule
\multicolumn{2}{l}{\textbf{Variables}}\\
\midrule
$t_i^a,t_i^d \geq 0$ & arrival/departure time at stop $ i \in \mathcal{I}$. \\
$\delta_i \in \{0,1\}$ & 1 if service at customer $ i \in \mathcal{C}$ starts outside its window.\\
$y_i \in \{0,1\}$ & 1 if the truck charges at stop $ i \in \mathcal{K}$.\\
$\tau_i^c \geq 0$ & charging duration at stop $ i \in \mathcal{K}$.\\
$e_i^a, e_i^d \geq 0$ & SOC at arrival/departure at CS at stop $ i \in \mathcal{I}$.\\
$\lambda_{i,r}^{a},\lambda_{i,r}^{d} \geq 0$ & PWL weight for encoding $e_i^a$, $e_i^d$, $i \in \mathcal{K}$, $r \in \mathcal{R}$.\\
$\mu_{i,r}^{a},\mu_{i,r}^{d} \in \{0,1\}$ & $1$ if $e_i^a$,$e_i^d$ lies in segment $r$ of the PWL curve, $i \in \mathcal{K}$, $r \in \mathcal{R}$.\\
$c_i^d \geq 0$ & accumulated consecutive driving time at arrival at stop $i$ (h).\\
$s_i^d \geq 0$ & accumulated shift driving time at arrival at stop $i$ (h).\\
$s_i^w \geq 0$ & accumulated shift working time at arrival at stop $i$ (h).\\
$h_i \geq 0$ & Elapsed time since end of last daily rest (h).\\
$x_i^{b_{45}}, x_i^{b_{15}}, x_i^{b_{30}} \in \{0,1\}$ & break type indicators at stop $i$. \\
$\phi_i \in \{0,1\}$&  1 if the first split break has been taken since the last driving reset.\\
$\tau_i^{b} \geq 0$& Break time beyond charging at stop $i$. \\
$\hat{\tau}_i^{b} \geq 0$& Break duration credited for HoS purposes. \\
$g_i \geq 0$ & Charging time credited as break at stop $i$. \\
$\sigma_i \in \{0,1\}$ & 1 if CS $i$ operates in sequential mode (no charging/break synchronization).\\
$\rho_i^{r_1}, \rho_i^{r_2} \in \{0,1\}$& rest type indicators at stop $i$.\\
$\tau_i^r \geq 0$& rest duration at stop $i$.\\
$v_i \in \{0,1\}$ & 1 if any activity (charge/break/rest) occurs at stop $i$.\\
$z_i, q_i \in \{0,1\}$ & extended shift flag; extension consumed at a rest.\\
\bottomrule
\multicolumn{2}{l}{{\footnotesize Abbreviations: HoS = Hours-of-service, CS = Charging station}}\\
\end{tabular}
\end{table}

\subsection{Sets and Parameters}
Let the route consist of a fixed sequence of stops indexed by $i \in \mathcal{I} = \{0, 1, \ldots, N\}$, where stop $0$ is the origin and stop $N$ is the destination. The set of customer stops is $\mathcal{C} \subseteq \mathcal{I}$, and the set of CS stops is $\mathcal{K} \subseteq \mathcal{I}$. Additionally, rest areas $\mathcal{L} \subseteq \mathcal{I}$ are place along the route and are treated as Charging station but with charging disabled. The set of non-customer stops $\mathcal S$ include both charging station and rest areas, such that  $\mathcal{C} \cap \mathcal{K} \cap \mathcal{L} = \emptyset$ and $\mathcal{K} \cup \mathcal{L} = \mathcal S$.  The model structure is similar to that of \citet{bai_rollout-based_2023}, adapted here to a single-vehicle fixed-route setting with multiple customer stops, as seen in Figure \ref{fig:routeExample}.
The nonlinear charging curve is approximated by a piecewise linear function defined by $R+1$ breakpoints, following \citet{montoya_electric_2017}.

\begin{figure}[htbp]
    \centering
    \fcolorbox{annblue}{white}{\parbox[c][2.6cm][c]{0.9\linewidth}{\centering
    {\color{annblue}\small\textsf{FIGURE NOT IN THE REPOSITORY --- \texttt{figures/routeExample.png}}}\\[3pt]
    {\color{annblue}\footnotesize Present on Overleaf; reproduced here as a placeholder only.\\ Your caption is kept below and is correct.}}}
    \caption{Structure of a benchmark corridor. Charging stations (triangles) are placed on a 60 km grid,  rest-area nodes (ticks) subdivide any leg longer than 25 km and allow for a break or rest but no charging, customers (diamonds) are drawn around cluster centers at 25\%, 50\% and 70\% of the corridor. The instance shown is a short-route case: 900 km, three customers.}
    \label{fig:routeExample}
\end{figure}

\subsection{Objective Function}
The objective minimizes the arrival time at the destination plus a fixed penalty $\beta$ for every customer whose service starts outside its window.
\begin{equation}
    \min \quad t_N^a + \beta \sum_{i \in \mathcal C}\delta_i
\end{equation}

\begin{acmt}
A referee at TR-E will ask why the objective is makespan rather than cost. The
answer is defensible --- on a fixed route with a fixed vehicle, driver time
dominates the marginal cost and asset utilisation scales with makespan --- but it
should be said here in one sentence, and the headline number converted once into
money in Section~\ref{sec:discussion}.
\end{acmt}

\subsection{Constraints}
In order to simplify the following constraints, we define additional variables as such: $\rho_i = \rho_i^{r1} + \rho_i^{r2} \in \{0,1\}$ is the rest indicator at a stop, $x_i = x_i^{b45} + x_i^{b15} + x_i^{b30} \in \{0,1\}$ is the break indicator at a stop, $r_i = x_i^{b45} + x_i^{b30} + \rho_i^{r1} + \rho_i^{r2} \in \{0,1\}$ is the reset indicator of consecutive drive time at a stop. We further write $o_i=t^d_i - t^a_i - \tau_i^r $ for the on-duty (non-rest) time spent at stop $i$, used in the worktime shift counter.

\subsubsection{Time propagation}
Constraint~\eqref{eq:time_prop} states that the arrival time at each stop equals the departure time from the previous stop plus the leg travel time. The departure time from a stop depends on the activities undertaken there. At a customer stop, constraint~\eqref{eq:depart_C} accounts for mandatory service time $S_i$ and any additional break or rest duration. At a CS stop, constraint~\eqref{eq:depart_K} additionally includes a queue time $Q_i y_i$, a charging duration $\tau_i^c$, a maneuver time $M_i^{stop}$ to get to/leave the CS in case we stop there for any activity, a maneuver time $M_i^{seq}$ if we first charge and then sequentially take a break or rest at the CS. Constraints~\eqref{eq:tw_soft1}--\eqref{eq:tw_soft2} activate the out-of-window indicator $\delta_i$ whenever service starts before the window opens or after it closes. Since service starts at arrival, the window applies to $t_i^a$.
\begin{equation}
    t_{i+1}^a = t_i^d + D_i, \qquad \, i \in \mathcal{I}\setminus\{N\} \label{eq:time_prop}
\end{equation}
\begin{equation}
    t_i^d = t_i^a + S_i + \tau_i^{b} + \tau_i^r, \qquad \, i \in \mathcal{C} \label{eq:depart_C}
\end{equation}
\begin{equation}
    t_i^d = t_i^a +  M_i^{stop} v_i + Q_i y_i + \tau_i^c + M_i^{seq} \sigma_i + \tau_i^{b} + \tau_i^r, \qquad \, i \in \mathcal{S} \label{eq:depart_K}
\end{equation}
\begin{equation}
    t_i^a \leq W_i^{f} + M \delta_i, \qquad \, i \in \mathcal{C} \label{eq:tw_soft1}
\end{equation}
\begin{equation}
    t_i^a \geq W_i^{a} - M \delta_i, \qquad \, i \in \mathcal{C} \label{eq:tw_soft2}
\end{equation}

\subsubsection{Battery state-of-charge}
Constraint~\eqref{eq:soc_prop} propagates the SOC between stops according to the energy consumed on each leg. Constraints~\eqref{eq:soc_nc}--\eqref{eq:soc_mono} enforce that charging can only occur at designated CS stops and can never decrease the SOC. Constraints~\eqref{eq:soc_lb}--\eqref{eq:soc_ub} keep the SOC within the feasible battery operating range at all times. Constraint~\eqref{eq:chg_act} ensures that a positive charging duration can only be incurred when the binary $y_i$ is active, and that in case of charging, we charge at least as much as the minimum charging time allowed.
\begin{equation}
    e_{i+1}^a = e_i^d - E_i, \qquad \, i \in \mathcal{I}\setminus\{N\} \label{eq:soc_prop}
\end{equation}
\begin{alignat}{2}
    e_i^d &= e_i^a, &\qquad& \, i \in \mathcal{C} \cup \mathcal L \label{eq:soc_nc}\\
    e_i^d &\geq e_i^a, &\qquad& \, i \in \mathcal{K} \label{eq:soc_mono}
\end{alignat}
\begin{alignat}{2}
    e_i^a &\geq E^{\min}, &\qquad& \, i \in \mathcal{I} \label{eq:soc_lb}\\
    e_i^d &\leq E^{\mathrm{cap}}, &\qquad& \, i \in \mathcal{I} \label{eq:soc_ub}
\end{alignat}
\begin{equation}
    \mathcal{T}_{\text{min}} y_i \leq \tau_i^c \leq \mathcal{T}_R\, y_i, \qquad \, i \in \mathcal{K} \label{eq:chg_act}
\end{equation}

\begin{acmt}
$E^{\min} = 0.2\,E^{cap}$ is a modelling choice (range anxiety plus battery
health), not a physical limit, and it removes 100\,kWh of usable capacity ---
roughly one 90\,km leg. Worth one justifying sentence with a citation in
Section~\ref{sec:exp}, since every stranding rate in
Table~\ref{tab:feasibility} is measured against this threshold rather than
against an empty battery.
\end{acmt}

\subsubsection{Piecewise-Linear Charging Function}
Following \citet{montoya_electric_2017}, we linearize the nonlinear charging function using a convex-combination formulation. Constraints~\eqref{eq:pwl_ea}--\eqref{eq:pwl_ed} express both the arrival and departure SOC at each CS stop as convex combinations of the $R+1$ PWL breakpoints. The charging duration is obtained in constraint~\eqref{eq:pwl_tauc} as the difference in cumulative charge time between the departure and arrival SOC levels. Constraints~\eqref{eq:pwl_ca}--\eqref{eq:pwl_cd} impose the convex-combination normalization condition on the arrival and departure weights. Constraints~\eqref{eq:sos2_sa}--\eqref{eq:sos2_sd} enforce that exactly one PWL segment is active for each SOC level. Finally, the SOS2 adjacency constraints~\eqref{eq:sos2_lo}--\eqref{eq:sos2_hi} restrict non-zero weights to at most two consecutive breakpoints, for both arrival ($\nu=a$) and departure ($\nu=d$) SOC levels.
\begin{align}
    e_i^a &= \sum_{r \in \mathcal{R}} \lambda_{i,r}^{a}\,\bar{E}_r, \qquad \, i \in \mathcal{K} \label{eq:pwl_ea}\\
    e_i^d &= \sum_{r \in \mathcal{R}} \lambda_{i,r}^{d}\,\bar{E}_r, \qquad \, i \in \mathcal{K} \label{eq:pwl_ed}
\end{align}
\begin{equation}
    \tau_i^c = \sum_{r \in \mathcal{R}} \lambda_{i,r}^{d}\,\mathcal{T}_r - \sum_{r \in \mathcal{R}} \lambda_{i,r}^{a}\,\mathcal{T}_r, \qquad \, i \in \mathcal{K} \label{eq:pwl_tauc}
\end{equation}
\begin{alignat}{2}
    \sum_{r \in \mathcal{R}} \lambda_{i,r}^{a} &= 1, &\qquad& \, i \in \mathcal{K} \label{eq:pwl_ca}\\
    \sum_{r \in \mathcal{R}} \lambda_{i,r}^{d} &= 1, &\qquad& \, i \in \mathcal{K} \label{eq:pwl_cd}
\end{alignat}
\begin{alignat}{2}
    \sum_{r=1}^{R} \mu_{i,r}^{a} &= 1, &\qquad& \, i \in \mathcal{K} \label{eq:sos2_sa}\\
    \sum_{r=1}^{R} \mu_{i,r}^{d} &= 1, &\qquad& \, i \in \mathcal{K} \label{eq:sos2_sd}
\end{alignat}
\begin{align}
    \lambda_{i,0}^{\nu}   &\leq \mu_{i,1}^{\nu}, \label{eq:sos2_lo}\\
    \lambda_{i,r}^{\nu}   &\leq \mu_{i,r}^{\nu} + \mu_{i,r+1}^{\nu}, \quad r = 1,\ldots,R-1, \label{eq:sos2_mid}\\
    \lambda_{i,R}^{\nu}   &\leq \mu_{i,R}^{\nu}, \label{eq:sos2_hi}
\end{align}
for all $i\in\mathcal{K}$ and $\nu\in\{a,d\}$.

\subsubsection{Break and Rest Decisions}
Constraint~\eqref{eq:one_act} ensures that at most one type of break or rest activity is declared at each stop.
\begin{equation}
    x_i    + \rho_i \leq 1, \qquad \, i \in \mathcal{I} \label{eq:one_act}
\end{equation}

At a CS stop, charging time may qualify toward HoS break requirements only when a break is declared and runs in parallel with charging ($x_i = 1$, $\sigma_i = 0$).
Let $g_i$ denote the credited charging time. Then constraints \eqref{eq:pi1}--\eqref{eq:pi3} force $g_i = \tau_i^c$ exactly in that case and $g_i = 0$ otherwise.
The effective break duration $\hat \tau_i^b$ is then the credited charging time plus any additional break time \eqref{eq:bhat_K}, and the additional break time alone at customer and rest-area stops \eqref{eq:bhat_C}.
\begin{equation}
    g_i \leq \tau_i^c \qquad \, i \in \mathcal{K} \label{eq:pi1}
\end{equation}
\begin{equation}
    g_i \leq \mathcal T_R (1- \sigma_i) \qquad \, i \in \mathcal{K} \label{eq:pi2}
\end{equation}
\begin{equation}
    g_i \leq \mathcal T_R x_i \qquad \, i \in \mathcal{K} \label{eq:pi21}
\end{equation}
\begin{equation}
    g_i \geq \tau_i^c - \mathcal T_R \sigma_i - \mathcal T_R (1-x_i) \qquad \, i \in \mathcal{K} \label{eq:pi3}
\end{equation}
\begin{alignat}{2}
    \hat{\tau}_i^b &= g_i+ \tau_i^b, &\quad& \, i \in \mathcal{K} \label{eq:bhat_K}\\
    \hat{\tau}_i^b &= \tau_i^b,            &\quad& \, i \in \mathcal{C} \cup \mathcal L \label{eq:bhat_C}
\end{alignat}

\begin{anew}
This block is the formulation's real contribution and it is stated only
algebraically. One worked sentence would carry it:
``A 50-minute charge at a station where a full break is declared in parallel
discharges that break at zero marginal time cost; a 30-minute charge does not, so
$\sigma_i$ is forced to one and the 45-minute break is appended after the charge
at a cost of $M^{seq}_i + 45$\,min. The whole diesel comparison of
Section~\ref{sec:diesel} is an accounting of this mechanism.''
\end{anew}

Constraints~\eqref{eq:brk45}--\eqref{eq:brk30} impose the minimum effective break duration for each break type. Similarly, constraints~\eqref{eq:rst1}--\eqref{eq:rst2} impose the minimum rest duration for each rest type. Constraints~\eqref{eq:brk_ub}--\eqref{eq:rst_ub} force the break and rest durations to zero when no corresponding activity is declared. Constraint~\eqref{eq:r2_lim} enforces that the reduced daily rest of $9$\,h may be used at most $\bar \rho$ times \DEL{per week}\FIX{per route}.
\begin{alignat}{2}
    \hat{\tau}_i^b &\geq T^{\mathrm{brk}}_{45}\, x_i^{b_{45}}, &\quad& \, i \in \mathcal{I} \label{eq:brk45}\\
    \hat{\tau}_i^b &\geq T^{\mathrm{brk}}_{15}\, x_i^{b_{15}}, &\quad& \, i \in \mathcal{I} \label{eq:brk15}\\
    \hat{\tau}_i^b &\geq T^{\mathrm{brk}}_{30}\, x_i^{b_{30}}, &\quad& \, i \in \mathcal{I} \label{eq:brk30}
\end{alignat}
\begin{alignat}{2}
    \tau_i^r &\geq T^{\mathrm{rst}}_{1}\, \rho_i^{r_1}, &\quad& \, i \in \mathcal{I} \label{eq:rst1}\\
    \tau_i^r &\geq T^{\mathrm{rst}}_{2}\, \rho_i^{r_2}, &\quad& \, i \in \mathcal{I} \label{eq:rst2}
\end{alignat}
\begin{alignat}{2}
    \tau_i^b &\leq M\, x_i,    &\quad& \, i \in \mathcal{I} \label{eq:brk_ub}\\
    \tau_i^r &\leq M\, \rho_i, &\quad& \, i \in \mathcal{I} \label{eq:rst_ub}
\end{alignat}
\begin{equation}
    \sum_{i \in \mathcal{I}} \rho_i^{r_2} \leq \bar{\rho} \label{eq:r2_lim}
\end{equation}

EU regulation allows the standard $45$\,min break to be replaced by a split break taken in two parts: a first part of at least $15$\,min followed by a second part of at least $30$\,min, in that strict order. The binary $\phi_i$ tracks whether the first part of a split break has been completed since the last consecutive-driving reset. Constraint~\eqref{eq:split_ord} ensures the second part can only be declared if the first part was taken at an earlier stop. Constraints~\eqref{eq:phi1}--\eqref{eq:phi3} propagate $\phi$ between stops. It is set to $1$ when a first split-break part is taken and no reset event has occurred, it is cleared to $0$ by any consecutive-driving reset (a full break, the completion of the second split-break part, or any rest).
\begin{equation}
    x_i^{b_{30}} \leq \phi_i, \qquad \, i \in \mathcal{I} \label{eq:split_ord}
\end{equation}
\begin{alignat}{3}
    \phi_{i+1} &\geq \phi_i + x_i^{b_{15}} - r_i, &\quad& \, i \in \mathcal{I}\setminus\{N\} \label{eq:phi1}\\
    \phi_{i+1} &\leq \phi_i + x_i^{b_{15}}, &\quad& \, i \in \mathcal{I}\setminus\{N\} \label{eq:phi2}\\
    \phi_{i+1} &\leq 1 - r_i, &\quad& \, i \in \mathcal{I}\setminus\{N\} \label{eq:phi3}
\end{alignat}

Constraint~\eqref{eq:anyactivity} flags any activity at a stop of $\mathcal S$. Sequential mode $\sigma_i$ is possible only when charging occurs \eqref{eq:sigma1}. It is forced whenever charging coexists with a rest \eqref{eq:sigma2}, since the truck must vacate the charging bay before an 11\,h or 9\,h rest, and it is forced whenever a break is declared that the charge does not fully cover \eqref{eq:sigma3}, since the driver must then take the break separately after charging. In \eqref{eq:sigma3}, $\underline b_i = T^{brk}_{45}x^{b_{45}}_i + T^{brk}_{15}x^{b_{15}}_i + T^{brk}_{30}x^{b_{30}}_i$ is the required minimum of the break declared at stop $i$; if the credited charge $g_i$ (equivalently, since $g_i = \tau^c_i$ under parallel mode, the charge itself) falls short of it, $\sigma_i$ is forced to 1, and constraint~\eqref{eq:pi2} then sets $g_i = 0$ so that the full break sits sequentially after the charge.
\begin{alignat}{2}
    \tfrac{1}{2}(y_i + \rho_i + x_i) \;\leq\; v_i &\;\leq\; y_i + \rho_i + x_i, &\qquad& i \in \mathcal{S} \label{eq:anyactivity}\\
    \sigma_i &\leq y_i, &\qquad& i \in \mathcal{K} \label{eq:sigma1}\\
    \sigma_i &\geq y_i + \rho_i - 1, &\qquad& i \in \mathcal{K} \label{eq:sigma2}\\
    \sigma_i &\geq \frac{\underline b_i - \tau^c_i}{T^{brk}_{45}} - (1 - y_i), &\qquad& i \in \mathcal{K} \label{eq:sigma3}
\end{alignat}

\subsubsection{Driving regulation}
The accumulated consecutive driving time $c_i^d$ increases with each driven leg and is reset to zero by any qualifying break or rest (i.e.\ whenever $r_i = 1$). Defining the auxiliary $l_i^1 = r_i\, c_i^d$ as the portion of $c_i^d$ that is zeroed out at stop $i$, constraint~\eqref{eq:cd_prop} expresses the propagation rule. Constraints~\eqref{eq:l1_ub1}--\eqref{eq:l1_lb} linearize the bilinear product $l_i^1 = r_i\, c_i^d$ via a standard big-$M$ reformulation. Constraint~\eqref{eq:cd_ub} enforces the regulatory upper bound on consecutive driving time.
\begin{equation}
    c_{i+1}^d = c_i^d + D_i - l_i^1, \qquad \, i \in \mathcal{I}\setminus\{N\} \label{eq:cd_prop}
\end{equation}
\begin{equation}
    c_i^d \leq T^{\mathrm{drv}}_{\mathrm{cons}}, \qquad \, i \in \mathcal{I} \label{eq:cd_ub}
\end{equation}

The accumulated shift driving time $s_i^d$ is reset only by a daily rest (a break alone is insufficient). Defining $l_i^2 = \rho_i\, s_i^d$, constraint~\eqref{eq:sd_prop} gives the propagation rule. Constraints~\eqref{eq:l2_ub1}--\eqref{eq:l2_lb} linearize $l_i^2 = \rho_i\, s_i^d$ with big-M $T^{drv}_{sh_2} = 10$\,h, the largest value $s^d_i$ can legitimately attain on an extended shift. Constraint~\eqref{eq:sd_ub} imposes the maximum shift driving time, conditional on the extension flag $z_i$. Constraint~\eqref{eq:z_prop} lets the flag persist within a shift and reset only after a rest. Constraints~\eqref{eq:q_def}--\eqref{eq:ext_budget} charge one unit of the \DEL{weekly} budget $\bar\epsilon$ per extended shift, the term $z_N$ covering an unfinished final shift.
\begin{equation}
    s_{i+1}^d = s_i^d + D_i - l_i^2, \qquad \, i \in \mathcal{I}\setminus\{N\} \label{eq:sd_prop}
\end{equation}
\begin{equation}
    s_i^d \leq T^{\mathrm{drv}}_{\mathrm{sh_1}} +(T^{\mathrm{drv}}_{\mathrm{sh_2}} - T^{\mathrm{drv}}_{\mathrm{sh_1}}) z_i , \qquad \, i \in \mathcal{I}. \label{eq:sd_ub}
\end{equation}
\begin{alignat}{2}
    z_{i+1} &\geq z_i - \rho_i, &\qquad& i \in \mathcal{I}\setminus\{N\} \label{eq:z_prop}\\
    q_i &\geq z_i + \rho_i - 1, &\qquad& i \in \mathcal{I} \label{eq:q_def}\\
    \sum_{i \in \mathcal I} q_i + z_N &\leq \bar\epsilon && \label{eq:ext_budget}
\end{alignat}

\subsubsection{Working regulation}
Working time comprises driving, customer service, CS queuing, maneuvering, and charging time not covered by a synchronized break. The work contribution of charging is simply $\tau^c_i - g_i$, which is linear, so no auxiliary variable is needed. The accumulated shift working time $s_i^w$ is reset only by a daily rest.By defining $l^4_i = \rho_i s^w_i$,  constraints~\eqref{eq:l4_ub1}--\eqref{eq:l4_lb} provide the linearization with big-M $T^{spr}_2 = 15$\,h (within a shift, working time is bounded only by the spread). Constraint~\eqref{eq:sw_prop} propagates $s^w_i$. The counter carries no daily cap: the daily-level bound on on-duty time is the shift spread \eqref{eq:h_prop}--\eqref{eq:h_ub} below\FIX{, and the counter itself is retained for reporting and for the ex-post verification of the Directive's working-time provisions}.
\begin{equation}
    s_{i+1}^w = s_i^w - l_i^4 + D_i
    + \begin{cases}
         M_{i+1}^{stop} v_{i+1} + Q_{i+1}\, y_{i+1} + (\tau^c_{i+1} - g_{i+1}) + M_{i+1}^{seq} \sigma_{i+1}  & \text{if } i+1 \in \mathcal{S}, \\
        S_{i+1}                      & \text{if } i+1 \in \mathcal{C},
    \end{cases}
    \quad \, i \in \mathcal{I}\setminus\{N\}
    \label{eq:sw_prop}
\end{equation}

\begin{acmt}
\textbf{Good --- the 60\,h weekly working-time constraint has been removed from
this version.} That was the most dangerous item in the previous draft: it was in
the paper but explicitly out of scope in \texttt{MILP.py}, so every reported
number had been produced by a different model from the one described. All that
remains is to delete $T^{wk}_{60}$ from Table~\ref{tab:notation-mini} (done
above) and the ``approach the weekly driving cap'' sentence in
Section~\ref{sec:instances}.
\end{acmt}

The shift spread $h_i$ measures the elapsed time since the end of the last daily rest. Since a daily rest is always the last activity performed at a stop, the elapsed time from the end of the previous rest to the start of a rest at stop $i$ is exactly $h_i + o_i$, where $o_i = t^d_i - t^a_i - \tau^r_i$ is the on-duty stop time. Constraint~\eqref{eq:h_prop} propagates the spread, with the reset $l^5_i = \rho_i (h_i + o_i)$ linearized in \eqref{eq:l5_1}--\eqref{eq:l5_3}. Constraint~\eqref{eq:spread_cap} bounds the pre-rest spread at $T^{spr}_1 = 13$\,h when the upcoming rest is regular and $T^{spr}_2 = 15$\,h when it is reduced, implementing the requirement that the daily rest be completed within 24\,h of the end of the previous one. Constraint~\eqref{eq:h_ub} is the global bound for shift length.
\begin{equation}
    h_{i+1} = h_i + o_i + D_i - l^5_i, \qquad i \in \mathcal I \setminus \{N\}
    \label{eq:h_prop}
\end{equation}
\begin{alignat}{2}
    h_i + o_i &\leq T^{spr}_1 + (T^{spr}_2 - T^{spr}_1)\,\rho^{r_2}_i + T^{spr}_2\,(1 - \rho_i), &\qquad& i \in \mathcal I \label{eq:spread_cap}\\
    h_i &\leq T^{spr}_2, &\qquad& i \in \mathcal I \label{eq:h_ub}
\end{alignat}

\begin{anew}
\textbf{Promote this to a remark.} The equivalence is a genuine, if small,
modelling contribution, and it is the one component of the formulation with no
counterpart in the reference model of Section~\ref{sec:validation}. Suggested:

\medskip
\noindent\textbf{Remark 1.} \emph{Enforcing \eqref{eq:spread_cap} is equivalent to
requiring that each daily rest be completed within 24\,h of the end of the
previous one.} A regular rest lasts 11\,h and a reduced rest 9\,h; since a rest is
always the terminal activity at its stop, the 24\,h requirement binds the on-duty
spread preceding it at $24-11 = 13$\,h and $24-9 = 15$\,h respectively. This
replaces a rolling-window constraint, which would need time-indexed variables, by
a single state variable propagated along the route.
\end{anew}

\subsubsection{Initial Conditions}
Constraints~\eqref{eq:init_t_e}--\eqref{eq:init_sw_phi} fix the state at the origin: the truck departs at a given time $t^0$, with initial SOC $E^0$, and all HoS accumulators and the split-break flag are initialized to zero. No break or rest activity is permitted at the origin or at the destination: $x_0 = \rho_0 = x_N = \rho_N = 0$.
\begin{alignat}{2}
    t_0^a &= t^0, &\quad e_0^a &= E^0, \label{eq:init_t_e}\\
    c_0^d &= 0,   &\quad s_0^d &= 0,   \label{eq:init_cd_sd}\\
    s_0^w &= 0,   &\quad \phi_0 &= 0.  \label{eq:init_sw_phi}
\end{alignat}

\begin{figure}[htbp]
    \centering
    \fcolorbox{annblue}{white}{\parbox[c][2.6cm][c]{0.9\linewidth}{\centering
    {\color{annblue}\small\textsf{FIGURE NOT IN THE REPOSITORY --- \texttt{figures/ExampleSolution3.png}}}\\[3pt]
    {\color{annblue}\footnotesize Present on Overleaf; reproduced here as a placeholder only.}}}
    \caption{A feasible schedule on a 950 km corridor with three customers. The four panels share a common time axis. \NEW{\emph{Route}: the stops used by the solution. \emph{Truck}: driving, parked, and the intervals during which the vehicle occupies a charging bay. \emph{SOC}: state of charge between $E^{\min}$ and $E^{cap}$. \emph{HoS}: the consecutive-driving and shift-driving accumulators against their limits. \emph{Driver}: the corresponding driver activity. The schedule illustrates the three regimes of \eqref{eq:pi1}--\eqref{eq:sigma3}: at 284\,km a 45-min break runs in parallel with the charge and is discharged at no time cost ($g_i = \tau^c_i$); at 497\,km a 9\,h reduced rest forces $\sigma_i = 1$, so the truck vacates the bay and the charge earns no credit; at 540\,km and 779\,km a 15-min and a 30-min charge discharge the two halves of a split break.}}
    \label{fig:solution}
\end{figure}

\begin{acmt}
This is the most communicative object in the paper and its caption is one line.
A reader who understands it does not need \eqref{eq:pi1}--\eqref{eq:sigma3} to
follow Section~\ref{sec:diesel}; a reader who does not will struggle throughout.
Consider also moving it \emph{before} the model, as a motivating figure.
\end{acmt}
